% GENERATED FILE. Edit canonical lessons, then run npm run generate:books.
% canonical-source-hash: fnv1a64:84143ce5e35da701
% canonical-lessons: PT-C324-conto, PT-C324-narrar, PT-C324-narrativa, PT-C324-narrador, PT-C324-enredo, PT-R324-revisao-317-324

\chapter{Era uma vez}
\label{ch:pt-era-uma-vez}
\hlchaptermodality{324}

\begin{chapteropening}
\textbf{By the end of this chapter:} \emph{I can talk about stories: tales, narrating, the narrative, the narrator and the plot.}

Complete the last lesson of chapter 324: I can talk about stories: tales, narrating, the narrative, the narrator and the plot.
\end{chapteropening}

\section[o conto]{o conto — the short story, the tale}

\label{lesson:PT-C324-conto}



\begin{quote}
Before the new one: say the Portuguese for the Easter bread, then the Portuguese for the sugared almonds.
\end{quote}

\subsection*{You'll want to know: o conto}
\textbf{o conto} — \textquotedblleft{}the short story, the tale\textquotedblright{}.

\begin{quote}
O avô contava-nos um conto antes de dormir.
\end{quote}

\begin{quote}
\emph{Grandfather used to tell us a story before bed.}
\end{quote}

Built on \emph{contar}, ‘to tell’: \emph{contar um conto}.

\begin{cousinweb}
\emph{conto} is built on \emph{contar} \textquotedblleft{}to count, to tell\textquotedblright{}, which is inherited from Latin \emph{computāre} \textquotedblleft{}to reckon\textquotedblright{} (Lewis \& Short): to tell events one by one, like counting.
\end{cousinweb}

\begin{grammarlens}[title={What you've built}]
One more word for telling, explaining and describing at B1.
\end{grammarlens}

\subsection*{Your turn}
\begin{itemize}
\raggedright
  \item \emph{Say it:} \emph{o conto}
  \item \emph{Say it:} \emph{o conto}, once more
  \item \emph{Say it:} a sentence of your own with \emph{o conto}
  \item \emph{Recall:} say the Portuguese for the Easter bread, then the Portuguese for the sugared almonds, then say \emph{o conto} again
\end{itemize}

\begin{tcolorbox}[breakable,colback=teal!4,colframe=teal!35!black,title={Before you move on}]
What does o conto mean? (\textquotedblleft{}The short story, the tale\textquotedblright{}.) Say it once more.
\end{tcolorbox}
\section[narrar]{narrar — to narrate, to tell}

\label{lesson:PT-C324-narrar}



\begin{quote}
Before the new one: say the Portuguese for the sugared almonds, then the Portuguese for the short story.
\end{quote}

\subsection*{You'll want to know: narrar}
\textbf{narrar} — \textquotedblleft{}to narrate, to tell\textquotedblright{}.

\begin{quote}
O livro narra a vida de um pescador.
\end{quote}

\begin{quote}
\emph{The book tells the life of a fisherman.}
\end{quote}

More formal than \emph{contar}.

\begin{cousinweb}
\emph{narrar} is a learned borrowing from Latin \emph{nārrāre} \textquotedblleft{}to tell, to relate\textquotedblright{} (Lewis \& Short).
\end{cousinweb}

\begin{grammarlens}[title={What you've built}]
One more word for telling, explaining and describing at B1.
\end{grammarlens}

\subsection*{Your turn}
\begin{itemize}
\raggedright
  \item \emph{Say it:} \emph{narrar}
  \item \emph{Say it:} \emph{narrar}, once more
  \item \emph{Say it:} a sentence of your own with \emph{narrar}
  \item \emph{Recall:} say the Portuguese for the sugared almonds, then the Portuguese for the short story, then say \emph{narrar} again
\end{itemize}

\begin{tcolorbox}[breakable,colback=teal!4,colframe=teal!35!black,title={Before you move on}]
What does narrar mean? (\textquotedblleft{}To narrate, to tell\textquotedblright{}.) Say it once more.
\end{tcolorbox}
\section[a narrativa]{a narrativa — the narrative, the story as told}

\label{lesson:PT-C324-narrativa}



\begin{quote}
Before the new one: say the Portuguese for the short story, then the Portuguese for to narrate.
\end{quote}

\subsection*{You'll want to know: a narrativa}
\textbf{a narrativa} — \textquotedblleft{}the narrative, the story as told\textquotedblright{}.

\begin{quote}
A narrativa começa numa aldeia pequena.
\end{quote}

\begin{quote}
\emph{The narrative begins in a small village.}
\end{quote}

Also an adjective: \emph{um texto narrativo}.

\begin{cousinweb}
\emph{narrativa} is built on \emph{narrativo} (EtymDB), from Latin \emph{nārrātīvus}, from \emph{nārrāre} \textquotedblleft{}to tell\textquotedblright{}.
\end{cousinweb}

\begin{grammarlens}[title={What you've built}]
One more word for telling, explaining and describing at B1.
\end{grammarlens}

\subsection*{Your turn}
\begin{itemize}
\raggedright
  \item \emph{Say it:} \emph{a narrativa}
  \item \emph{Say it:} \emph{a narrativa}, once more
  \item \emph{Say it:} a sentence of your own with \emph{a narrativa}
  \item \emph{Recall:} say the Portuguese for the short story, then the Portuguese for to narrate, then say \emph{a narrativa} again
\end{itemize}

\begin{tcolorbox}[breakable,colback=teal!4,colframe=teal!35!black,title={Before you move on}]
What does a narrativa mean? (\textquotedblleft{}The narrative, the story as told\textquotedblright{}.) Say it once more.
\end{tcolorbox}
\section[o narrador]{o narrador — the narrator}

\label{lesson:PT-C324-narrador}



\begin{quote}
Before the new one: say the Portuguese for to narrate, then the Portuguese for the narrative.
\end{quote}

\subsection*{You'll want to know: o narrador}
\textbf{o narrador} — \textquotedblleft{}the narrator\textquotedblright{}.

\begin{quote}
O narrador é um menino de dez anos.
\end{quote}

\begin{quote}
\emph{The narrator is a ten-year-old boy.}
\end{quote}

A woman is \emph{a narradora}.

\begin{cousinweb}
\emph{narrador} is a learned borrowing from Latin \emph{nārrātor} \textquotedblleft{}a teller, a narrator\textquotedblright{} (EtymDB), from \emph{nārrāre} \textquotedblleft{}to tell\textquotedblright{}.
\end{cousinweb}

\begin{grammarlens}[title={What you've built}]
One more word for telling, explaining and describing at B1.
\end{grammarlens}

\subsection*{Your turn}
\begin{itemize}
\raggedright
  \item \emph{Say it:} \emph{o narrador}
  \item \emph{Say it:} \emph{o narrador}, once more
  \item \emph{Say it:} a sentence of your own with \emph{o narrador}
  \item \emph{Recall:} say the Portuguese for to narrate, then the Portuguese for the narrative, then say \emph{o narrador} again
\end{itemize}

\begin{tcolorbox}[breakable,colback=teal!4,colframe=teal!35!black,title={Before you move on}]
What does o narrador mean? (\textquotedblleft{}The narrator\textquotedblright{}.) Say it once more.
\end{tcolorbox}
\section[o enredo]{o enredo — the plot}

\label{lesson:PT-C324-enredo}



\begin{quote}
Before the new one: say the Portuguese for the narrative, then the Portuguese for the narrator.
\end{quote}

\subsection*{You'll want to know: o enredo}
\textbf{o enredo} — \textquotedblleft{}the plot\textquotedblright{}.

\begin{quote}
O enredo do filme é muito complicado.
\end{quote}

\begin{quote}
\emph{The plot of the film is very complicated.}
\end{quote}

Built on \emph{enredar}, ‘to tangle up’.

\begin{cousinweb}
\emph{enredo} is built on \emph{enredar} \textquotedblleft{}to catch in a net, to entangle\textquotedblright{}, from \emph{rede} \textquotedblleft{}a net\textquotedblright{}, which is inherited from Latin \emph{rēte} \textquotedblleft{}a net\textquotedblright{} (Lewis \& Short).
\end{cousinweb}

\begin{grammarlens}[title={What you've built}]
One more word for telling, explaining and describing at B1.
\end{grammarlens}

\subsection*{Your turn}
\begin{itemize}
\raggedright
  \item \emph{Say it:} \emph{o enredo}
  \item \emph{Say it:} \emph{o enredo}, once more
  \item \emph{Say it:} a sentence of your own with \emph{o enredo}
  \item \emph{Recall:} say the Portuguese for the narrative, then the Portuguese for the narrator, then say \emph{o enredo} again
\end{itemize}

\begin{tcolorbox}[breakable,colback=teal!4,colframe=teal!35!black,title={Before you move on}]
What does o enredo mean? (\textquotedblleft{}The plot\textquotedblright{}.) Say it once more.
\end{tcolorbox}
\section[(dialogue)]{Revisão dos capítulos 317–324}

\label{lesson:PT-R324-revisao-317-324}



\begin{quote}
Say the Portuguese for to retire and the pensioner.
\end{quote}

\subsection*{Your turn}
Say these aloud:

\begin{itemize}
\raggedright
  \item to retire, the pensioner, to emigrate, emigration, the emigrant
  \item to celebrate, to commemorate, the commemoration, the toast, to toast
  \item the saint, the march, the bush basil, the sardine feast, the street dance
  \item the custom, the ritual, the ceremony, the pilgrim, the pilgrimage
  \item the carnival, the parade, the mask, the disguise, to disguise
  \item the fireworks, the rocket, the bonfire, to decorate, the decoration
  \item the Christmas Eve dinner, the nativity scene, the New Year celebration, the Easter bread, the sugared almonds
  \item the short story, to narrate, the narrative, the narrator, the plot
\end{itemize}

\begin{tcolorbox}[breakable,colback=teal!4,colframe=teal!35!black,title={Before you move on}]
To retire? (\textbf{aposentar}.) The pensioner? (\textbf{o pensionista}.) To emigrate? (\textbf{emigrar}.) Emigration? (\textbf{a emigração}.) The emigrant? (\textbf{o emigrante}.) To celebrate? (\textbf{festejar}.) To commemorate? (\textbf{comemorar}.) The commemoration? (\textbf{a comemoração}.) The toast? (\textbf{o brinde}.) To toast? (\textbf{brindar}.) The saint? (\textbf{o santo}.) The march? (\textbf{a marcha}.) The bush basil? (\textbf{o manjerico}.) The sardine feast? (\textbf{a sardinhada}.) The street dance? (\textbf{o bailarico}.) The custom? (\textbf{o costume}.) The ritual? (\textbf{o ritual}.) The ceremony? (\textbf{a cerimónia}.) The pilgrim? (\textbf{o peregrino}.) The pilgrimage? (\textbf{a peregrinação}.) The carnival? (\textbf{o carnaval}.) The parade? (\textbf{o desfile}.) The mask? (\textbf{a máscara}.) The disguise? (\textbf{o disfarce}.) To disguise? (\textbf{disfarçar}.) The fireworks? (\textbf{o fogo de artifício}.) The rocket? (\textbf{o foguete}.) The bonfire? (\textbf{a fogueira}.) To decorate? (\textbf{enfeitar}.) The decoration? (\textbf{o enfeite}.) The Christmas Eve dinner? (\textbf{a consoada}.) The nativity scene? (\textbf{o presépio}.) The New Year celebration? (\textbf{a passagem de ano}.) The Easter bread? (\textbf{o folar}.) The sugared almonds? (\textbf{as amêndoas}.) The short story? (\textbf{o conto}.) To narrate? (\textbf{narrar}.) The narrative? (\textbf{a narrativa}.) The narrator? (\textbf{o narrador}.) The plot? (\textbf{o enredo}.)
\end{tcolorbox}
